\documentclass[11pt]{article}
\usepackage{amsmath,amssymb,amsthm,bm}
\usepackage[margin=1in]{geometry}
\usepackage{hyperref}
\usepackage{xcolor}
\usepackage{enumitem}
\usepackage{booktabs}

% -- Csurös-style notation (CM09 SI / Csurös 2021 arXiv) ------------------
\newcommand{\R}{\mathbb{R}}
\newcommand{\N}{\mathbb{N}}
\newcommand{\Z}{\mathbb{Z}}
\newcommand{\E}{\mathbb{E}}
\newcommand{\PROB}{\mathbb{P}}
\newcommand{\Tree}{\mathcal{T}}
\newcommand{\Leaves}{\mathcal{L}}
\newcommand{\pa}{\mathrm{pa}}
\newcommand{\ch}{\mathrm{ch}}
\newcommand{\troot}{R}
\newcommand{\Prof}{\boldsymbol{\xi}}
\newcommand{\rate}{\kappa}             % gain shape
\newcommand{\dup}{\lambda}             % duplication rate
\newcommand{\loss}{\mu}                % loss rate (=1 by convention)
\newcommand{\edgelen}{t}               % edge length
\newcommand{\survp}{\tilde p}          % survival probability
\newcommand{\survq}{\tilde q}          % Pólya per-event parameter
\newcommand{\Ll}{\mathcal{L}}
\newcommand{\PGF}{\mathsf{G}}
\DeclareMathOperator{\Pol}{Polya}
\DeclareMathOperator{\Poi}{Pois}
\DeclareMathOperator{\NB}{NegBin}
\DeclareMathOperator{\Surv}{Surv}

\newtheorem{theorem}{Theorem}
\newtheorem{proposition}[theorem]{Proposition}
\newtheorem{lemma}[theorem]{Lemma}
\newtheorem{corollary}[theorem]{Corollary}
\theoremstyle{definition}
\newtheorem{definition}[theorem]{Definition}
\newtheorem{remark}[theorem]{Remark}

\title{Boundary identifiability of the gain--loss--duplication likelihood:\\
       the Pólya--Poisson asymptote and what it means for\\
       ancestral genome inference}
\author{\texttt{recount} project\\
        \textit{notation following CM09 SI and Cs\H{u}r\"os 2021 arXiv}}
\date{\today}

\begin{document}
\maketitle

\begin{center}
\fbox{\parbox{0.9\textwidth}{\small\textbf{Provenance and disclaimer.}
This document was drafted by Claude (Anthropic's LLM, model
\texttt{claude-opus-4-7[1m]}) in collaboration with the human author,
as part of an interactive Claude Code session
(\url{https://claude.com/claude-code}). The mathematical derivations
follow standard arguments from Numerical Recipes 10.7, L\'evy
continuity for probability generating functions, and the survival
map of the CM09 supplement. They have been spot-checked numerically
against the analytical native C implementation but not externally
peer-reviewed. The arc269 numerical findings are reproducible from
\texttt{validation/} scripts (see
\texttt{validation/outputs/README.md}) but reflect the optimiser
trajectories of a particular software stack (\texttt{native\_bfgs}
at the time of writing) and may shift with optimiser choice or
precision. Treat as ``LLM-drafted, human-checked, FD-validated and
bit-perfectly reproduced against Cs\H{u}r\"os' Java on the focal
subsets'' rather than peer-reviewed manuscript content.}}
\end{center}

\begin{abstract}
We give a rigorous account of a boundary pathology in the
gain--loss--duplication (GLD) phylogenetic likelihood: there exist
sequences of interior parameters along which the log-likelihood
$\Ll(\Theta)$ increases without ever attaining its supremum at any
interior point. The mechanism is the Pólya--Poisson PGF convergence
$\Pol(\rate,\survq) \rightharpoonup \Poi(\rate\survq/(1-\survq))$
(Lemma~\ref{lem:pgf}), which one of the per-edge offspring kernels
can be driven to approach as $\dup_{v_\star}\to\infty$ at one
internal node $v_\star$. We characterise one sufficient
parameter-space path realising the limit
(Lemma~\ref{lem:pgf}, Remark~\ref{rem:suffroute}), establish
non-attainment along it (Theorem~\ref{thm:unbounded}), prove a
quantitative gap result that any finite log-rate budget leaves an
$\epsilon$-suboptimal interior MLE
(Proposition~\ref{prop:gap-quant}), and prove that a weakly
informative log-Normal prior on the rates renders the MAP objective
coercive with global maximum in the interior
(Theorem~\ref{thm:map-coercive}). The arc269 archaeal dataset
exhibits this asymptote in practice
(\S\ref{sec:numerical}); on a single internal node the optimiser
finds $\dup_{v_\star} \approx 5.6\times10^9$ giving an LL
improvement of $\sim 20{,}000$ nat over any in-basin interior point.
\end{abstract}

% =================================================================
\section{Biological motivation: what the asymptote means}
\label{sec:bio-motivation}
% =================================================================

In ancestral genome reconstruction, every internal node $v$ of the
species tree carries a hypothesis about how families arose,
duplicated, and were lost along the edge that leads to it. A pure
Pólya/negative-binomial offspring distribution gives \emph{lineage-by-lineage}
heritable expansion: copies that survive can themselves duplicate, and
the variance grows in step with the mean. A pure Poisson offspring
distribution gives instead a \emph{lineage-independent immigration}
process at the node: families just \emph{appear}, with no link to
what was already there. These are biologically different stories.

In the GLD parameterisation the Pólya--Poisson limit is the
``$\dup \to \infty$'' boundary: very fast duplication at one node,
combined with a very short observable survival fraction, makes the
per-edge kernel look like an immigration burst.  In other words, when
the optimiser pushes $\dup_{v_\star}$ to infinity at some node
$v_\star$, the model is saying ``\emph{at $v_\star$ many families
spontaneously materialised, in a way that is statistically
indistinguishable from a Poisson immigration source}.''

For most biological datasets this is a much better fit than the
finite-$\dup$ Pólya alternative --- the over-dispersion of
finite-$\dup$ Pólya is hard to make compatible with empirical
copy-number variance, but an unconstrained immigration ``vent'' at
one node can absorb whatever excess copies are needed.  Hence the
optimiser climbs the boundary and the inferred ancestral genome
size is then determined by whatever finite log-rate the optimiser
happens to stall at, \emph{not} by any biological information.

\paragraph{What this document does.}
\S\ref{sec:notation} sets up notation. \S\ref{sec:pgf} proves the
Pólya--Poisson PGF limit (Lemma~\ref{lem:pgf}) and identifies
a sufficient parameter-space path realising it
(Remark~\ref{rem:suffroute}).  \S\ref{sec:nonattain} establishes
non-attainment along that path (Theorem~\ref{thm:unbounded}) and a
quantitative gap result (Proposition~\ref{prop:gap-quant}).
\S\ref{sec:map} proves coercivity of the MAP objective under the
log-Normal prior (Theorem~\ref{thm:map-coercive}).
\S\ref{sec:numerical} documents the asymptote on arc269.
\S\ref{sec:biology} discusses biological implications and
method-selection recommendations.

% =================================================================
\section{Setup and notation}
\label{sec:notation}
% =================================================================

We follow Cs\H{u}r\"os \& Mikl\'os 2009 SI (henceforth \textbf{CM09})
and Cs\H{u}r\"os 2021 arXiv:2107.11440 (henceforth \textbf{Cs21}); the
reader is assumed familiar with the GLD setup, this section fixes
notation only.

\paragraph{Tree, profiles, rates.}
Let $\Tree = (V, E, \troot)$ be a rooted tree with $N=|V|$ nodes,
edges directed from the root $\troot$ to the leaves
$\Leaves \subset V$. For $v \neq \troot$ write $\pa(v)\in V$ for the
parent, $\ch(v) \subset V$ for the children. On the incoming edge of
each non-root node we attach rates
\begin{equation}
\rate_v \;\ge\; 0, \qquad
\loss_v > 0, \qquad
\dup_v \;\ge\; 0, \qquad
\edgelen_v \;\in\; (0, \infty),
\label{eq:rates}
\end{equation}
where $\rate_v$ is the gain (Pólya shape) parameter, $\loss_v$ the
per-lineage loss rate (Cs21 convention $\loss_v \equiv 1$),
$\dup_v$ the per-lineage duplication rate, $\edgelen_v$ the edge
time, and the root convention $\edgelen_\troot = +\infty$. A
\emph{gene-family profile} is a vector
$\Prof \in \Z_{\ge 0}^\Leaves$ of integer copy counts at each leaf;
the data is $\{\Prof_f\}_{f=1}^F$. The free parameter is
$\Theta = (\rate_v, \dup_v, \edgelen_v)_{v\ne \troot} \cup
\{\rate_\troot\}$, taking values in the \emph{open} set
\begin{equation}
\Theta^\circ
\;=\;
(\R_{>0})^N_\rate \;\times\; (\R_{>0})^{N-1}_\dup \;\times\;
(\R_{>0})^{N-1}_\edgelen.
\label{eq:openparam}
\end{equation}

\paragraph{The conserved-copy decomposition.}
Cs21 decomposes the transition $\xi_{\pa(v)} = n \to \xi_v = m$ on
edge $\pa(v)\to v$ via the hidden variable
$\eta_v = $ ``\emph{number of $n$'s lineages that have at least one
descendant at $v$}'':
\begin{align}
\PROB\bigl[\eta_v = s \,\bigl|\, \xi_{\pa(v)} = n\bigr]
&= \binom{n}{s}\,(1-\survp_v)^s\,\survp_v^{n-s},
\qquad 0 \le s \le n,
\label{eq:eta-given-xi}\\
\PROB\bigl[\xi_v = m \,\bigl|\, \eta_v = s\bigr]
&= \binom{\rate_v + m - 1}{m - s}\,
   (1-\survq_v)^{\rate_v + s}\,\survq_v^{m - s},
\qquad m \ge s.
\label{eq:xi-given-eta}
\end{align}
The two per-edge probabilities $\survp_v, \survq_v$ are determined by
$(\loss_v, \dup_v, \edgelen_v)$ via the survival map.  Following the
CM09 SI we use the convention that $\survp_v$ is the probability that
a single lineage at the top of the edge has \emph{no} surviving
descendant at the bottom (the parent of $\eta_v$); thus
$1-\survp_v$ in \eqref{eq:eta-given-xi} is the per-lineage
\emph{survival} probability. The closed form is
\begin{equation}
\survp_v
=
\frac{\loss_v\bigl(e^{(\loss_v - \dup_v)\edgelen_v} - 1\bigr)}
     {\loss_v\,e^{(\loss_v - \dup_v)\edgelen_v} - \dup_v},
\qquad
\survq_v
=
\frac{\dup_v\bigl(e^{(\loss_v - \dup_v)\edgelen_v} - 1\bigr)}
     {\loss_v\,e^{(\loss_v - \dup_v)\edgelen_v} - \dup_v},
\label{eq:surv-explicit}
\end{equation}
for $\dup_v > 0$, with the Poisson collapse $\survp_v = e^{-\loss_v\edgelen_v}$,
$\survq_v = 0$ when $\dup_v = 0$. The map
$(\loss_v, \dup_v, \edgelen_v) \mapsto (\survp_v, \survq_v)$ is
smooth on the interior of \eqref{eq:openparam}.

\paragraph{Per-edge transition kernel.}
Composing \eqref{eq:eta-given-xi} and \eqref{eq:xi-given-eta}:
\begin{equation}
T_v(m \mid n)
\;:=\;
\sum_{s=0}^{\min(n,m)}
   \binom{n}{s}\,(1-\survp_v)^s\,\survp_v^{n-s}
   \cdot \Pol\bigl(m - s \,\big|\, \rate_v + s,\, \survq_v\bigr),
\label{eq:transition}
\end{equation}
where the Pólya (negative-binomial) PMF is
\begin{equation}
\Pol(k \mid \alpha, q)
\;=\;
\frac{\Gamma(\alpha + k)}{\Gamma(\alpha)\, k!}\,
q^k (1 - q)^\alpha,
\qquad k \in \Z_{\ge 0},
\label{eq:polya-pmf}
\end{equation}
i.e.\ the Pólya distribution with shape $\alpha$ and tail $q$.
Per-family likelihoods $L_f(\Theta) = \PROB[\Prof_f \mid \Theta]$
are assembled by the inside--outside recursion of CM09; the total
log-likelihood is $\Ll(\Theta) = \sum_{f=1}^F \log L_f(\Theta)$.

\paragraph{Goal.}
Characterise
\begin{equation}
\Ll^\star \;:=\; \sup_{\Theta\in\Theta^\circ}\Ll(\Theta).
\label{eq:supgoal}
\end{equation}

% =================================================================
\section{The Pólya--Poisson PGF convergence}
\label{sec:pgf}
% =================================================================

The point of the next lemma is to identify a parameter-space path
along which the per-edge offspring distribution converges to a
Poisson. We use probability generating functions and L\'evy
continuity (which says that pointwise convergence of PGFs on $[0,1]$
is equivalent to weak convergence of integer-valued distributions).

The PGF of $\Pol(\alpha, q)$ is, for $|s| < q^{-1}$,
\begin{equation}
\PGF_{\alpha, q}(s) \;=\; \E[\,s^X\,]
\;=\;
\Bigl(\tfrac{1-q}{1 - q\,s}\Bigr)^{\!\alpha},
\label{eq:polya-pgf}
\end{equation}
and the PGF of $\Poi(\Lambda)$ is
$\PGF^{\mathrm{Pois}}_\Lambda(s) = e^{\Lambda(s-1)}$.

\begin{lemma}[Pointwise Pólya--Poisson convergence]\label{lem:pgf}
Fix $\Lambda > 0$. Along any sequence
$(\alpha_n, q_n) \in \R_{>0}\times(0,1)$ with
$\alpha_n \to \infty$, $q_n \to 0$, and
$\alpha_n q_n / (1-q_n) \to \Lambda$,
\begin{equation}
\PGF_{\alpha_n, q_n}(s) \;\longrightarrow\; e^{\Lambda(s-1)}
\qquad \forall s \in [0, 1].
\label{eq:pgf-conv}
\end{equation}
Consequently $\Pol(\alpha_n, q_n) \rightharpoonup \Poi(\Lambda)$
weakly on $\Z_{\ge 0}$.
\end{lemma}
\begin{proof}
From \eqref{eq:polya-pgf}
\begin{equation*}
\log\PGF_{\alpha_n, q_n}(s)
\;=\;
\alpha_n\bigl[\log(1-q_n) - \log(1 - q_n s)\bigr].
\end{equation*}
For $s \in [0, 1]$ and $q_n \in (0,1)$ both arguments lie in $(0,1)$,
so the Taylor expansion at $0$ is valid:
\begin{equation*}
\log(1 - q_n) = -q_n - \tfrac{q_n^2}{2} + O(q_n^3),
\qquad
\log(1 - q_n s) = -q_n s - \tfrac{(q_n s)^2}{2} + O(q_n^3),
\end{equation*}
with $O$ constant uniform in $s\in[0,1]$. Substituting,
\begin{equation*}
\log\PGF_{\alpha_n, q_n}(s)
\;=\;
\alpha_n q_n (s - 1)
\;+\; \alpha_n q_n^2 \cdot \tfrac{1}{2}(s^2 - 1)
\;+\; O(\alpha_n q_n^3).
\end{equation*}
By hypothesis $\alpha_n q_n \to \Lambda(1 - q_n) \to \Lambda$
(since $q_n\to 0$), and
$\alpha_n q_n^2 = (\alpha_n q_n) \cdot q_n \to \Lambda \cdot 0 = 0$,
and similarly $\alpha_n q_n^3 \to 0$. Hence the right-hand side
tends to $\Lambda(s-1)$ pointwise, giving \eqref{eq:pgf-conv}.
Pointwise convergence of PGFs on $[0,1]$ to a continuous limit at
$s = 1$ implies weak convergence of distributions on $\Z_{\ge 0}$
(L\'evy continuity for PGFs of integer-valued random variables).
$\square$
\end{proof}

\begin{remark}[A sufficient parameter-space path to Lemma~\ref{lem:pgf}]
\label{rem:suffroute}
Lemma~\ref{lem:pgf}'s hypothesis is met along $\dup_v \to \infty$
at fixed $(\loss_v, \edgelen_v)$ \emph{provided} $\rate_v$ is
co-rescaled. Write $z := e^{(\loss_v - \dup_v)\edgelen_v}$. As
$\dup_v \to \infty$ this decays exponentially to $0$. From
\eqref{eq:surv-explicit},
\begin{align*}
\survq_v &= \tfrac{\dup_v(1 - z)}{\dup_v - \loss_v z}
       \;\sim\; 1 - z\;\longrightarrow\; 1,
\\
1 - \survq_v &\sim z = e^{-(\dup_v - \loss_v)\edgelen_v}
\quad(\text{exponentially small}),
\\
\survp_v &\;\sim\; \tfrac{\loss_v z}{\dup_v - \loss_v z}
       \;\sim\; \tfrac{\loss_v}{\dup_v}
       \;\longrightarrow\; 0
\quad(\text{algebraically}).
\end{align*}
For fixed $\rate_v$, the would-be Poisson mean is
$\rate_v \survq_v / (1 - \survq_v) \sim \rate_v / z
= \rate_v e^{(\dup_v - \loss_v)\edgelen_v} \to \infty$,
so Lemma~\ref{lem:pgf}'s hypothesis is \emph{not} met with $\rate_v$
fixed. To meet it one must shrink $\rate_v$ along the path with
$\rate^\star_v \cdot z^{-1} \to \Lambda_v$, i.e.\
$\rate^\star_v \sim \Lambda_v z \to 0$ at the same exponential rate
at which $1 - \survq_v$ vanishes.
This co-rescaling is one \emph{sufficient} path on
$(\rate_v, \dup_v) \in \Theta^\circ$ realising the per-edge Poisson
limit. We will see (Remark~\ref{rem:empirical-path}) that the
optimiser empirically realises a \emph{different}, broader-tailed
path, but the unboundedness phenomenon is the same.
\end{remark}

\begin{proposition}[Convergence of per-edge transition kernels]
\label{prop:edge-conv}
For fixed bounded $m, n \in \Z_{\ge 0}$, the per-edge transition
$T_v(m \mid n)$ in \eqref{eq:transition} satisfies, along the
sufficient path of Remark~\ref{rem:suffroute},
\begin{equation}
T_{v,n}(m \mid n) \;\longrightarrow\;
T^\star_v(m \mid n)
\;:=\;
\sum_{s=0}^{\min(n,m)} \binom{n}{s} (1-\survp^\star)^s (\survp^\star)^{n-s}\,
\Poi\bigl(m - s \,\big|\, \Lambda_v\bigr),
\label{eq:edge-limit}
\end{equation}
with $\survp^\star = \lim_n\survp_{v,n}$ and $\Lambda_v$ the limiting
Poisson mean of Lemma~\ref{lem:pgf}.
\end{proposition}
\begin{proof}
The binomial factor $\binom{n}{s}(1-\survp)^s\survp^{n-s}$ is
continuous in $\survp$. The Pólya factor
$\Pol(m - s \mid \rate + s, \survq)$ converges to
$\Poi(m - s \mid \Lambda)$ by Lemma~\ref{lem:pgf}, with shape
$\rate + s$ finite. Each summand of \eqref{eq:transition}
converges, and the sum is finite ($\le \min(n,m) + 1$ terms);
limits commute with finite sums. $\square$
\end{proof}

\begin{corollary}[Likelihood limit along the sufficient path]
\label{cor:LL-sup-attain}
Along the Remark~\ref{rem:suffroute} sequence
$\Theta_n \to (\Theta^\circ)^{\dup_{v_\star} = \infty}$, the
per-family likelihood limit $L^\star_f := \lim_n L_f(\Theta_n) \in (0,1]$
exists for every $f$, and $\log L_f$ extends continuously from
$\Theta^\circ$ to the single boundary point of the path.
\end{corollary}
\begin{proof}
Apply Proposition~\ref{prop:edge-conv} at each edge and propagate
through the finite-dimensional inside--outside recursion of CM09,
which is a polynomial assembly of edge kernels and is therefore
continuous in each. $\square$
\end{proof}

% =================================================================
\section{Non-attainment of the supremum}
\label{sec:nonattain}
% =================================================================

The next theorem turns Corollary~\ref{cor:LL-sup-attain} into a
non-attainment statement: the LL strictly increases along the
asymptote but never reaches its limit at any finite parameter.

\begin{theorem}[Non-attainment along the sufficient path]
\label{thm:unbounded}
Let $\Theta_n$ be the Remark~\ref{rem:suffroute} sequence at boundary
node $v_\star$, and let $L^\star_f := \lim_n L_f(\Theta_n)$ as in
Corollary~\ref{cor:LL-sup-attain}. Suppose there exists a family
$f_0$ for which
\begin{equation}
L_{f_0}(\Theta_n) \;<\; L^\star_{f_0}
\qquad \text{for all } n \text{ sufficiently large.}
\label{eq:strict-suff}
\end{equation}
Then the limiting value
$\Ll^\star_{\mathrm{path}} := \sum_f \log L^\star_f$ is finite,
equals the supremum of $\Ll$ along the path, and is not attained at
any point of the path: $\Ll(\Theta_n) < \Ll^\star_{\mathrm{path}}$
for every such $n$, with strict approach as $n\to\infty$.
\end{theorem}
\begin{proof}
By Corollary~\ref{cor:LL-sup-attain},
$\log L_f(\Theta_n) \to \log L^\star_f$ along the path for every $f$,
with $L^\star_f \in (0, 1]$, so the sum
$\Ll^\star_{\mathrm{path}} = \sum_f \log L^\star_f$ is finite (each
summand is $\le 0$ and the sum is over a finite index $f=1,\ldots,F$).
For $n$ large enough,
\begin{align*}
\Ll(\Theta_n) - \Ll^\star_{\mathrm{path}}
&\;=\; \sum_{f \neq f_0}\bigl[\log L_f(\Theta_n) - \log L^\star_f\bigr]
\;+\; \bigl[\log L_{f_0}(\Theta_n) - \log L^\star_{f_0}\bigr].
\end{align*}
The first sum tends to $0$, and the second is strictly negative by
\eqref{eq:strict-suff}. Hence
$\Ll(\Theta_n) < \Ll^\star_{\mathrm{path}}$ for all such $n$, and
$\Ll(\Theta_n) \uparrow \Ll^\star_{\mathrm{path}}$ from below.
$\square$
\end{proof}

\begin{remark}[Genericity of \eqref{eq:strict-suff}]
\label{rem:genericity}
Condition \eqref{eq:strict-suff} fails only when every family's
likelihood is already maximised by an interior point and the boundary
strictly worsens it. For real gene-family data this is rarely the
case: typically at least one family has a sub-Pólya empirical
variance at some clade and benefits from the Poisson asymptote at
the corresponding internal node.
\end{remark}

\begin{remark}[The empirical optimiser path is broader]
\label{rem:empirical-path}
On arc269 (\S\ref{sec:numerical}) the optimiser finds
$\dup_{v_\star} \approx 5.6\times10^9$ at one internal node, so
$1 - \survq_{v_\star}$ is of order $10^{-9}$ in $\edgelen_{v_\star}$.
But the corresponding $\rate_{v_\star} \approx 1.76\times 10^3$ ---
\emph{not} the exponentially small $\rate^\star_{v_\star}$ required
by Remark~\ref{rem:suffroute} (where one would need
$\rate_{v_\star} \sim z \approx 10^{-9}$). The empirical path
therefore has
$\rate_{v_\star} \survq_{v_\star}/(1 - \survq_{v_\star}) \to \infty$;
its per-edge kernel sends offspring mass to $+\infty$, not to a
proper Poisson limit. The Pólya PGF $\PGF_{\alpha_n, q_n}(s) \to 0$
for every $s\in[0,1)$ while $\PGF_{\alpha_n, q_n}(1) = 1$ --- so
the offspring distribution escapes to infinity.

\emph{Why does the LL still climb?} At $v_\star$ many families
require many copies to be ``emitted'' at that node, and the GLD
parameters elsewhere cannot supply them; pushing both
$\dup_{v_\star}$ and $\rate_{v_\star}$ large gives a very
heavy-tailed Pólya offspring with both mean and variance diverging,
and the joint LL improves until line-search precision stalls.
A precise analytical characterisation of \emph{this} broader
asymptote is beyond our scope; Lemma~\ref{lem:pgf} establishes that
\emph{one} mechanism gives unboundedness, and the empirical path
exhibits the same phenomenon within a wider class.
\end{remark}

\begin{proposition}[Quantitative gap]\label{prop:gap-quant}
Let $v_\star$ be any node. For every $\epsilon > 0$ there exists a
finite log-rate budget $B_\epsilon < \infty$ such that
\begin{equation}
\sup_{\substack{\Theta \in \Theta^\circ\\\log\dup_{v_\star} \le B_\epsilon}}
\Ll(\Theta)
\;\ge\;
\Ll^\star_{\mathrm{path}} - \epsilon,
\label{eq:gap}
\end{equation}
and the supremum is attained at a finite interior point.
\end{proposition}
\begin{proof}
By Corollary~\ref{cor:LL-sup-attain}, $\Ll(\Theta_n) \to \Ll^\star_{\mathrm{path}}$
along the Remark~\ref{rem:suffroute} sequence. Choose $n_0$ with
$\Ll(\Theta_{n_0}) > \Ll^\star_{\mathrm{path}} - \epsilon$ and set
$B_\epsilon = \log\dup_{v_\star, n_0}$.  The point $\Theta_{n_0}$ is
interior and witnesses \eqref{eq:gap}.  Since the constrained
problem ``$\log\dup_{v_\star} \le B_\epsilon$'' is the
intersection of $\Theta^\circ$ with a closed half-space, its
attainment requires only continuity of $\Ll$ and compactness of the
super-level set, which follow from the smoothness of the survival
map. $\square$
\end{proof}

\begin{remark}[What Theorem~\ref{thm:unbounded} does \emph{not} claim]
The theorem establishes non-attainment along the specific sufficient
asymptote of Remark~\ref{rem:suffroute}. It does not assert that
the global supremum
$\sup_{\Theta\in\Theta^\circ}\Ll$ equals $\Ll^\star_{\mathrm{path}}$:
a separate interior critical point with higher likelihood is not
ruled out. Numerically no such interior critical point has been
found on arc269 across our multistart sweeps
(\S\ref{sec:numerical}); the empirical optimum is always asymptote.
Theorem~\ref{thm:unbounded} is therefore conservative ---
non-attainment in the only direction the optimiser actually probes.
\end{remark}

% =================================================================
\section{Resolution: a coercive MAP objective}
\label{sec:map}
% =================================================================

The Pólya--Poisson asymptote is a property of the GLD likelihood
considered on $\Theta^\circ$ alone. Adding a weakly informative
log-Normal prior on the rates removes the pathology by making the
objective coercive on log-rate coordinates.

\begin{definition}[Weakly informative log-Normal prior]
For each non-root node $v$ and the root gain, define independently
\begin{equation}
\log\rate_v \sim \mathcal{N}(m_\rate, \sigma^2),
\qquad
\log\dup_v \sim \mathcal{N}(m_\dup, \sigma^2),
\qquad
\log\edgelen_v \sim \mathcal{N}(m_\edgelen, \sigma^2),
\label{eq:lognormal}
\end{equation}
with default centres $m_\rate = \log 0.1$, $m_\dup = \log 0.5$,
$m_\edgelen = \log 1$ (biologically reasonable defaults), and
$\sigma = 1$.
\end{definition}

The MAP objective is
\begin{equation}
\Ll^{\mathrm{MAP}}(\Theta)
\;:=\;
\Ll(\Theta) + \log p_\pi(\Theta)
\;=\;
\Ll(\Theta) \;-\; \tfrac{1}{2\sigma^2}\!\!\sum_{v\ne\troot}
   \bigl[(\log\rate_v - m_\rate)^2
       + (\log\dup_v - m_\dup)^2
       + (\log\edgelen_v - m_\edgelen)^2\bigr]
\;+\; \mathrm{const.}
\label{eq:mapobj}
\end{equation}

\begin{theorem}[Coercivity of the MAP objective]
\label{thm:map-coercive}
For any $\sigma < \infty$, the MAP objective \eqref{eq:mapobj} is
coercive on $\Theta^\circ$ in log-rate coordinates:
$\Ll^{\mathrm{MAP}}(\Theta_n) \to -\infty$ whenever any one of
$|\log\rate_{v,n}|, |\log\dup_{v,n}|, |\log\edgelen_{v,n}|$
tends to $+\infty$. The supremum of $\Ll^{\mathrm{MAP}}$ on
$\Theta^\circ$ is attained at an interior point.
\end{theorem}
\begin{proof}
For each family $f$, $L_f(\Theta) \in (0,1]$, so
$\log L_f(\Theta) \le 0$. Summing over the finitely many families,
$\Ll(\Theta) \le 0$ uniformly on $\Theta^\circ$, hence bounded
above.

The prior term in \eqref{eq:mapobj} is the quadratic form
$-\tfrac{1}{2\sigma^2}\|\log\Theta - m\|_2^2 + \mathrm{const}$
with full-rank precision $1/\sigma^2 \cdot I$, so it tends to
$-\infty$ at quadratic rate as any one of $|\log\rate_v|,
|\log\dup_v|, |\log\edgelen_v|$ goes to infinity. A bounded
likelihood plus a quadratic prior that tends to $-\infty$ gives
$\Ll^{\mathrm{MAP}}(\Theta_n)\to -\infty$, which is coercivity.

Coercivity together with continuity of $\Ll^{\mathrm{MAP}}$ (the
likelihood is smooth on $\Theta^\circ$ from the survival map; the
prior is a smooth quadratic) implies every super-level set
$\{\Theta : \Ll^{\mathrm{MAP}}(\Theta) \ge c\}$, for $c$ inside the
attained range, is non-empty, bounded by coercivity, hence compact
(in the log-rate coordinates) on which a continuous function attains
its supremum. $\square$
\end{proof}

\begin{remark}[Calibration of $\sigma = 1$]
With $\sigma = 1$ the 99\,\% prior credible interval for each rate
is $[m \cdot e^{-2.58}, m \cdot e^{2.58}]$, i.e.\ roughly two orders
of magnitude wide. The quadratic prior cost at
$\log\dup_v = \log(5.6\times10^9) - m_\dup \approx 22.5$ is
$22.5^2/2 \approx 253$ nat per such node, comfortably outweighing
the $\sim 20{,}000$-nat asymptote gain at a single node when
counted across the full prior; in practice the MAP optimum lands at
$\max_v\dup_v \approx 43$, $\max_v\rate_v \approx 0.12$ ---
biologically interpretable rates, far from any boundary.
\end{remark}

% =================================================================
\section{Numerical evidence: arc269}
\label{sec:numerical}
% =================================================================

\paragraph{Dataset.}
The 269-leaf archaeal phylogeny of Cs26 (\texttt{tree269-baker-alti-gtdb};
$F = 90{,}243$ families, $N = 537$ nodes, filter $\Omega_{\min} = 1$)
is the largest dataset bundled with the paper. We report fits using
the consolidated driver \texttt{validation/ml.py} with clip
$B = 50$ (rates in $[2\times10^{-22},\, 5\times10^{21}]$).

\paragraph{Multiple basins, all on the asymptote.}
Two distinct ML optimiser trajectories are observed.

\begin{itemize}[leftmargin=*]
\item \textbf{Simple-ML basin} (single warm-restart BFGS from default
init): $\Ll = -1{,}129{,}653$, $\bar\xi_\troot = 8{,}664$ root copies,
$3{,}533$ root families, copies/family $2.45$.  Rate distribution:
$\max\rate_v = 1{,}540$, $\max\dup_v = 7.6\times10^9$.
\item \textbf{Polish basin} (multistart + lognormal polish at
$\sigma = 0.10$): $\Ll = -1{,}109{,}266$
($+20{,}387$ nat over the simple-ML basin),
$\bar\xi_\troot = 8{,}963$ root copies,
$3{,}608$ root families, copies/family $2.48$. Rate distribution:
$\max\rate_v = 1{,}763$, $\max\dup_v = 5.63\times10^9$ (a single
internal node carries the entire boundary load; the second-largest
$\dup_v$ in the rate vector is $124$, eight orders of magnitude
below).
\end{itemize}
Both basins are sequences along the asymptote of
Remark~\ref{rem:empirical-path} at one node $v_\star$; the polish
basin has simply climbed further up the same asymptote
(cf.\ Proposition~\ref{prop:gap-quant} with smaller $\epsilon$).

\paragraph{Independent confirmation from the focal subsets.}
The same phenomenon is visible on the smaller subset clades
(D80, P75, E114, ED194), but softer: Cs26's published ML rates report
$\max\rate_v \approx 10^{14}$ on these subsets (Java GUI inspection
of the bundled \texttt{.countxml.gz}). At the subset scale this
drives only a $\sim 20\%$ spread in $\bar\xi_\troot$ between
alternative boundary climbs; at the arc269 scale (with $\sim 20\times$
more families and a larger tree), the same mechanism inflates
$\bar\xi_\troot$ by $2$--$3\times$.

\paragraph{MAP $\sigma=1$ fits land in the interior.}
On the four focal subsets, MAP with $\sigma=1$ lands within
$\pm 300$ nat of the published Cs26 ML values and within $\pm 20\%$
on $\bar\xi_\troot$, so the prior does not distort the
non-boundary-pathological subsets. On arc269, the MAP from default
uniform initialisation converges to raw $\Ll = -1{,}117{,}480$
(\,$-\log P_{\mathrm{MAP}} = 1{,}120{,}615$ including the
$-3{,}135$-nat prior cost), with $\bar\xi_\troot = 5{,}906$ root
copies, $3{,}108$ root families, copies/family $1.90$,
$\max\dup_v = 43$, $\max\rate_v = 0.12$ (interior, far from any
boundary). A 15-start lognormal-perturbation bootstrap discovers an
even deeper basin at $-\log P = 1{,}098{,}463$, raw
$\Ll = -1{,}096{,}512$, $\sim 22{,}000$ nat below the cold-start
basin --- showing that the MAP objective has multiple local maxima
even at $\sigma = 1$ and that multistart is required to locate the
global mode.

\begin{center}
\begin{tabular}{lrrrrr}
\toprule
Method & $\Ll$ & $\bar\xi_\troot$ & root fam.\ & cp/fam & $\max\dup_v$ \\
\midrule
Cold-start ML ($B=50$) & $-1{,}129{,}653$ & $8{,}664$ & $3{,}533$ & $2.45$ & $7.6\times10^9$ \\
Multistart-polish ML ($B=50$) & $-1{,}109{,}266$ & $8{,}963$ & $3{,}608$ & $2.48$ & $5.6\times10^9$ \\
MAP $\sigma=1$ (cold start) & $-1{,}117{,}480^{\dagger}$ & $5{,}906$ & $3{,}108$ & $1.90$ & $43$ \\
\textbf{MAP $\sigma=1$ (15-start global)} & $\mathbf{-1{,}096{,}512^{\dagger}}$ & $\mathbf{5{,}426}$ & $\mathbf{3{,}070}$ & $\mathbf{1.77}$ & $\mathbf{28}$ \\
\bottomrule
\end{tabular}
\end{center}
{\small $^{\dagger}$ pure log-likelihood at the MAP estimate, not
the regularised objective.}

% =================================================================
\section{Biological consequences and method-selection guidance}
\label{sec:biology}
% =================================================================

\begin{enumerate}[leftmargin=*]
\item \textbf{ML alone is not enough for ancestral reconstruction at
depth.} The two ML basins on arc269 differ only in how far along
the asymptote the optimiser climbed; they do not represent distinct
biological inferences. The biological quantity of interest, the
root genome size $\bar\xi_\troot$, is sensitive to the (essentially
arbitrary) stopping criterion of the optimiser.
\item \textbf{The MAP $\sigma=1$ 15-start global estimate} (root
copies $5{,}426$, root families $3{,}070$, copies/family $1.77$) is
our recommended LACA quote. The cold-start MAP basin is itself an
interior local optimum $\sim 22{,}000$ nat above the global ---
showing that even with the prior, multistart is required.
\item \textbf{Cs26's published focal-subset fits are also on the
asymptote.} With $\max\rate_v\approx 10^{14}$ on D80--ED194
(log-rate $\sim 32$, well below the $B=50$ clip but well above any
biologically interpretable range), every published reconstruction
is technically a boundary artefact; the small magnitude on the
4-subset fits keeps the practical impact below $\pm 20\%$ in
$\bar\xi_\troot$, but the underlying behaviour is the same as on
arc269.
\item \textbf{The boundary issue is a model feature, not a data
feature.} Any GLD reconstruction at depth without modelling
horizontal gene transfer (HGT) will exhibit the same asymptote ---
the optimiser is asking for a Poisson immigration source, and
without HGT in the model the only place to put one is at a Pólya
boundary.
\item \textbf{ALE-style models including HGT may shift LACA
estimates downward}, since families currently absorbed into the
``inflated ancestor + asymptote'' will be explained by horizontal
transfer events. The MAP $\sigma=1$ quote here is the cleanest
GLD-only baseline against which to measure that shift.
\end{enumerate}

% =================================================================
\appendix
\section{The transition-kernel limit at $\survp_v = 0$}
% =================================================================

For completeness we record the form of $T^\star_v$
(Proposition~\ref{prop:edge-conv}) when in addition
$\survp_v^\star = 0$ (parent-survival collapses): only the $s=0$ term
of the binomial survives, giving
\begin{equation}
T^\star_v(m \mid n) \;=\; \Poi(m \mid \Lambda_v),
\label{eq:edge-collapsed}
\end{equation}
\emph{independent of $n$}. The parent count is forgotten, so the
subtree below $v$ contributes a constant per-family Poisson factor.
This is the degenerate limit at which the inferred ancestor becomes
literally a Poisson immigration source.

% =================================================================
\section{Why a finite log-rate clip alone is a less elegant fix}
\label{app:clip}
% =================================================================

A finite log-rate clip $B<\infty$ is equivalent to a uniform improper
prior on $[m_a - B, m_a + B]$ in log-rate space (and $0$ outside).
This prior is non-smooth (a hard wall at $|\log\cdot| = B$), so the
clipped MLE can sit on the boundary of the clip set itself with
$\dup_v = e^B$.  With $B = 50$ in our numerics the clip is rarely
attained (the optimiser stalls at $\log\dup_{v_\star} \approx 22$
due to line-search precision), so the clipped problem behaves as
unbounded and inherits all the asymptote pathology characterised in
Theorem~\ref{thm:unbounded}--Proposition~\ref{prop:gap-quant}.

The MAP prior \eqref{eq:lognormal} achieves the same boundedness
through a \emph{smooth} quadratic cost. Both objectives are coercive,
but only the smooth one yields a unique optimum that is robust to
optimiser choice and stopping criterion.

\end{document}
